% Introduction to Part VI. Unnumbered text: no \label except the \phantomsection target.
\phantomsection\label{intro:VI}
Part~VI proves Ba\l aban's variational theorem on the torus for every number $k$ of averaging steps, with constants that do not depend on $k$. Theorem~\ref{thm:variational} solves the problem relative to a regular reference, but its regularity bound for the minimizer, $\eps+C_6(\eps+\beta)$, depends on the reference. In an induction over $k$ the reference at step $k$ is the solution at step $k-1$, and the change of units multiplies its regularity parameter by a power of $L$ (Lemma~\ref{lem:units}). Without an improvement the constants would grow geometrically in $k$. Part~VI supplies the improvement on the torus, without the nested domains of \cite[Sect.~F]{BalVar}.

\paragraph{What this Part proves.} The main result is Theorem~\ref{thm:balaban}, which with Corollary~\ref{cor:regularity} is Theorem~\ref{thm:J}. There are $a_1$, $B_3$ and $n_0$ such that, for all $k\ge0$, all tori of side $n\ge n_0$ and every real coarse configuration $V$ with curvature at most $\eps_1\le a_1$: the action on $\{U:U^{(k)}=V\}$ has exactly one critical orbit among configurations of small curvature and current; it consists of the minimizers there; the background fields $U_k(V)$ have curvature and current at most $B_3\eps_1$ and are $(M,B_3\eps_1)$-regular; and the minimal orbit of $rV$ is the image of that of $V$ for $r\in\mathfrak S_k$. The constants depend on $L$, $G$, $a$, $\alpha$ and $M$, and not on $k$, $n$ or $\eta$. Corollary~\ref{cor:regularity} shows that every real configuration with curvature and current at most $\theta\le a_2$ is $(M,B_5\theta)$-regular, so the hypothesis of Part~III holds on Ba\l aban's small-field class. Corollary~\ref{cor:analytic} gives analyticity of the background map around a background field.

\paragraph{How.} The reference at step $k$ is the solution of the $(k-1)$-step problem for the trivial extension of $V$, which has exactly the prescribed averages (Lemma~\ref{lem:trivial-extension}), and Theorem~\ref{thm:variational} applies to it. The local improvement (Lemma~\ref{lem:improvement}) restores the regularity bound: near a point $y$, the reference is replaced by a configuration that is trivial near $y$ in a local gauge and equals the minimizer far away (Definition~\ref{def:splice}, Lemma~\ref{lem:splice}). Relative to it, the first-order formula of Theorem~\ref{thm:variational}(c) shows that the minimizer near $y$ is determined by the coarse field in the gauge of the coordinate-ordered contours from $y$ (Lemma~\ref{lem:adapted-gauge}), which the curvature of $V$ controls. The errors are exponentially small in the splicing radius, or quadratic (Section~\ref{sec:local-gauges}, Proposition~\ref{prop:localized}). This gives $\eps_0'=C_{\rm imp}\eps_1+\frac12\eps_0$, and the induction closes. Gauges on cubes larger than the regularity scale come from Lemma~\ref{lem:large-cubes}, with bounds that grow polynomially in the size.

\paragraph{Imports.} Theorem~\ref{thm:variational} and its Euler--Lagrange equation (Proposition~\ref{prop:euler-lagrange}); from Part~IV, the regularity of averages and the averages in a local gauge (Proposition~\ref{prop:avg-regular}, Lemma~\ref{lem:avg-local-gauge}). One condition on $a_1$ in the proof of Theorem~\ref{thm:balaban}, $B_3a_1\le c'_3$ with $c'_3$ of Section~\ref{sec:AB-proof}, serves only Part~XI (proof of Theorem~\ref{thm:ol-chart}). No other proof in Part~VI uses a later result.

\paragraph{Exports.} The background fields $U_k(V)$ of Theorem~\ref{thm:balaban}, their regularity and their symmetry are used in Parts~VII and IX--XI. Corollary~\ref{cor:regularity} goes to Parts~VIII--X, Corollary~\ref{cor:analytic} to Parts~VII and IX--XI, and Lemma~\ref{lem:large-cubes} to Parts~VII--X. The coordinate-ordered contours and the splicing of Section~\ref{sec:local-gauges} are used again for the curvature bound of Proposition~\ref{prop:sharp-curvature}. Lemmas~\ref{lem:patching}, \ref{lem:adapted-gauge} and Proposition~\ref{prop:localized} serve only this Part. Linear bounds on cubes of size of order $\eps_1^{-1}$, as in \cite[(9), (10)]{BalVar}, are not proved; Theorem~\ref{thm:balaban} and Lemma~\ref{lem:large-cubes} give polynomial ones (Remark~\ref{rem:large-cubes}).

\paragraph{Details for this Part.} The companion has no chapter on Part~VI; its parameter ledger lists $a_1$, $B_3$, $n_0$ and the other constants of Part~VI that Parts~VII--XI use.
